% GENERATED FILE. Edit canonical lessons, then run npm run generate:books.
% canonical-source-hash: fnv1a64:09e598ca5aafe9c3
% canonical-lessons: TE-C59-wind, TE-C59-mist, TE-C59-sand, TE-C59-mud, TE-C59-fire

\chapter{Weather and Ground}
\label{ch:te-weather-ground}
\hlchaptermodality{59}

\begin{chapteropening}
\textbf{By the end of this chapter:} \emph{I can name the weather and what is underfoot in Telugu, and use it as a reason to go home.}

Complete TE-C59-fire and demonstrate the chapter promise.
\end{chapteropening}

\section[gāli]{\te{గాలి} (gāli) — wind, air}

\label{lesson:TE-C59-wind}



\begin{quote}
Before the new one: what did \emph{antē} mean?
\end{quote}

\subsection*{You'll want to know: \te{గాలి}}
\textbf{\te{గాలి}} (\emph{gāli}) — \textquotedblleft{}wind, air\textquotedblright{}.

\te{గాలి} is inherited, and it covers both the wind that blows and the air that is simply there. Telugu does not split them.

It reaches further than weather. \te{గాలి} is breath, and it is rumour — \te{గాలి} \te{వార్త} is a story with nothing holding it up, blown in from somewhere. Air, breath and hearsay under one word is a chain a great many languages make, and it usually starts with breathing.

A \te{గాలి} getting up in the evening is the ordinary reason given for cutting a visit short.

\begin{grammarlens}[title={What you've built}]
The first reason to start saying goodbye.
\end{grammarlens}

\subsection*{Your turn}
Say these aloud:

\begin{itemize}
\raggedright
  \item \emph{gāli}
  \item it once more, slowly
  \item \emph{gāli}, then \emph{veḷḷostānu}, and use the one as a reason for the other
\end{itemize}

\begin{tcolorbox}[breakable,colback=teal!4,colframe=teal!35!black,title={Before you move on}]
What does \te{గాలి} mean? (\textquotedblleft{}wind, air\textquotedblright{}.) Say it once more.
\end{tcolorbox}
\section[mañcu]{\te{మంచు} (mañcu) — mist, dew, snow}

\label{lesson:TE-C59-mist}



\begin{quote}
Before the new one: what did \emph{gāli} mean?
\end{quote}

\subsection*{You'll want to know: \te{మంచు}}
\textbf{\te{మంచు}} (\emph{mañcu}) — \textquotedblleft{}mist, dew, snow\textquotedblright{}.

\te{మంచు} is the cold wet thing in the air: mist on a winter morning, dew on the ground, and snow where there is any. One word for all three states, which suits a country that meets the third of them rarely.

This is worth sitting with. A language draws its lines where its speakers need them. Telugu splits its cooking vessels three ways and lumps mist, dew and snow together; a language spoken where snow falls would do the opposite.

\te{మంచు} \te{కురుస్తోంది} is \textquotedblleft{}it is misting\textquotedblright{} on the coast and \textquotedblleft{}it is snowing\textquotedblright{} in the hills, with no change to the word.

\begin{grammarlens}[title={What you've built}]
Two: what moves, and what settles.
\end{grammarlens}

\subsection*{Your turn}
Say these aloud:

\begin{itemize}
\raggedright
  \item \emph{mañcu}
  \item it once more, slowly
  \item \emph{mañcu}, then \emph{gāli}, and say which of the two you can see
\end{itemize}

\begin{tcolorbox}[breakable,colback=teal!4,colframe=teal!35!black,title={Before you move on}]
What does \te{మంచు} mean? (\textquotedblleft{}mist, dew, snow\textquotedblright{}.) Say it once more.
\end{tcolorbox}
\section[isuka]{\te{ఇసుక} (isuka) — sand}

\label{lesson:TE-C59-sand}



\begin{quote}
Before the new one: what did \emph{mañcu} mean?
\end{quote}

\subsection*{You'll want to know: \te{ఇసుక}}
\textbf{\te{ఇసుక}} (\emph{isuka}) — \textquotedblleft{}sand\textquotedblright{}.

\te{ఇసుక} is inherited and unremarkable, and that is the point: it is the ground beside every \te{నది} in the Telugu country.

The sand of a riverbed is where cattle are watered, where lorries come at night to dig, and where a village argument about who owns the river tends to start. A plain word can carry a lot of local weather.

\te{ఇసుకలో} \te{నీళ్ళు}, water in sand, is what you say about effort that soaks away.

\begin{grammarlens}[title={What you've built}]
Three.
\end{grammarlens}

\subsection*{Your turn}
Say these aloud:

\begin{itemize}
\raggedright
  \item \emph{isuka}
  \item it once more, slowly
  \item \emph{isuka}, then \emph{nadi}, and say where you would find one beside the other
\end{itemize}

\begin{tcolorbox}[breakable,colback=teal!4,colframe=teal!35!black,title={Before you move on}]
What does \te{ఇసుక} mean? (\textquotedblleft{}sand\textquotedblright{}.) Say it once more.
\end{tcolorbox}
\section[burada]{\te{బురద} (burada) — mud}

\label{lesson:TE-C59-mud}



\begin{quote}
Before the new one: what did \emph{isuka} mean?
\end{quote}

\subsection*{You'll want to know: \te{బురద}}
\textbf{\te{బురద}} (\emph{burada}) — \textquotedblleft{}mud\textquotedblright{}.

\te{బురద} is mud — what \te{ఇసుక} becomes when the rain has been at it.

It does the same second job that mud does in English. \te{బురద} \te{చల్లు}, to throw mud, is to slander somebody, and the picture is exact enough that you can guess the phrase before you are told it.

\te{బురద} is also the honest reason a visit ends early in the wet months. The \te{బాట} between the \te{గ్రామం} and the \te{నది} turns to it.

\begin{grammarlens}[title={What you've built}]
Four.
\end{grammarlens}

\subsection*{Your turn}
Say these aloud:

\begin{itemize}
\raggedright
  \item \emph{burada}
  \item it once more, slowly
  \item \emph{burada}, then \emph{isuka}, and say what turns one into the other
\end{itemize}

\begin{tcolorbox}[breakable,colback=teal!4,colframe=teal!35!black,title={Before you move on}]
What does \te{బురద} mean? (\textquotedblleft{}mud\textquotedblright{}.) Say it once more.
\end{tcolorbox}
\section[nippu]{\te{నిప్పు} (nippu) — fire, a live coal}

\label{lesson:TE-C59-fire}



\begin{quote}
Before the new one: what did \emph{burada} mean?
\end{quote}

\subsection*{You'll want to know: \te{నిప్పు}}
\textbf{\te{నిప్పు}} (\emph{nippu}) — \textquotedblleft{}fire, a live coal\textquotedblright{}.

\te{నిప్పు} is inherited, and it is the fire you can point at: the live coal, the burning thing in the hearth. \te{అగ్ని}, borrowed from Sanskrit, is fire as an idea and as a god, the one that carries offerings.

So Telugu keeps a working fire and a ceremonial one, in the same way it keeps a \te{పొద్దు} and a \te{సూర్యుడు}. That split has turned up so often in these chapters that you can start expecting it: the old word does the work, the borrowed one does the ceremony.

\te{నిప్పు} is what a house is banked down around at night, and putting it out is the last thing done before sleep. A good place to say \te{వెళ్ళొస్తాను}.

\begin{grammarlens}[title={What you've built}]
Five: \te{గాలి}, \te{మంచు}, \te{ఇసుక}, \te{బురద}, \te{నిప్పు}. Weather enough to take your leave on.
\end{grammarlens}

\subsection*{Your turn}
Say these aloud:

\begin{itemize}
\raggedright
  \item \emph{nippu}
  \item it once more, slowly
  \item all five in order, then say \emph{gāli ekkuva}, and take your leave with \emph{veḷḷostānu}
\end{itemize}

\begin{tcolorbox}[breakable,colback=teal!4,colframe=teal!35!black,title={Before you move on}]
What does \te{నిప్పు} mean? (\textquotedblleft{}fire, a live coal\textquotedblright{}.) Say it once more.
\end{tcolorbox}
